\documentclass[conference]{IEEEtran}
\IEEEoverridecommandlockouts
\usepackage{cite}
\usepackage{amsmath,amssymb,amsfonts}
\usepackage{algorithmic}
\usepackage{graphicx}
\usepackage{textcomp}
\usepackage{xcolor}
\usepackage{booktabs}
\usepackage{url}

\begin{document}

\title{SimQC -- Weekly Progress Report\\}
\author{
    \IEEEauthorblockN{Aoutul Nabi Purna\IEEEauthorrefmark{1}, Shishir Dhar\IEEEauthorrefmark{2}, Shihab Sharar\IEEEauthorrefmark{3}, and Md. Iftikhar Alam Hridik\IEEEauthorrefmark{4}}
    \IEEEauthorblockA{Department of Electrical and Computer Engineering, North South University, Dhaka, Bangladesh}
    \IEEEauthorblockA{Emails: \{\IEEEauthorrefmark{1}aoutul.purna.241, \IEEEauthorrefmark{2}shishir.dhar.241, \IEEEauthorrefmark{3}shihab.sharar.241, \IEEEauthorrefmark{4}iftikhar.hridik.241\}@northsouth.edu}
}

\maketitle

\begin{abstract}
This report outlines Sprint 2 of the SimQC quantum simulator, covering the state-vector engine core, the matrix-free pairwise gate-application kernel, the gate/parameter/registry layer, the circuit model with subcircuit inlining and flattening, the canonical JSON serializer, and the differential test harness validating the $O(2^n)$ kernel against the $O(4^n)$ Kronecker reference. The sprint translates the frozen conventions from Sprint 1 into a working zero-allocation simulation engine in JavaScript, with differential validation against the naive reference implementation for circuits up to six qubits.
\end{abstract}

\section{Progress Summary}
During this sprint, we implemented the simulation core specified in the frozen conventions from Sprint 1[cite: 1]. Completed deliverables include:

\begin{enumerate}
    \item A \texttt{StateVector} class exposing a clean facade (\texttt{get}/\texttt{set}, \texttt{probability}, \texttt{probabilities}, \texttt{normSquared}, \texttt{clone}, \texttt{renormalize}) while storing amplitudes in two flat \texttt{Float64Array} buffers (\texttt{re}, \texttt{im}), initialized to the zero state $\alpha_0 = 1$, all other amplitudes zero[cite: 1]. Internal buffer references \texttt{\_re} and \texttt{\_im} are exposed strictly for kernel access, preserving encapsulation for external consumers[cite: 1].

    \item A zero-allocation pairwise kernel in \texttt{kernel.js}. The \texttt{apply1QubitOp} function uses stride-$2^t$ block iteration over bitwise index pairs $(i, i \oplus (1 \ll t))$, with bitwise control masks \texttt{(i \& controlMask) === controlMask} and anti-control filtering. The inner loop unrolls the canonical $2\times2$ complex multiply-add on the flat 8-double tuple $g = [g_{00}^r, g_{00}^i, g_{01}^r, g_{01}^i, g_{10}^r, g_{10}^i, g_{11}^r, g_{11}^i]$, operating exclusively on primitive doubles with zero heap allocation[cite: 1]. A separate \texttt{applySwap} function implements SWAP/CNOT permutations via direct amplitude exchange without floating-point multiplication[cite: 1].

    \item A gate layer comprising \texttt{basis.js}, \texttt{param.js}, and \texttt{registry.js}. \texttt{basis.js} defines static gates ($I, X, Y, Z, H, S, T$) using Sprint 1's \texttt{Matrix} class for readability, then flattens each to an 8-tuple via \texttt{matrixToFlat8} at module initialization. \texttt{param.js} implements parameterized rotations ($R_x, R_y, R_z$) as pure functions of $\theta$ returning flat tuples directly. \texttt{registry.js} exposes an immutable map resolving string identifiers (\texttt{"H"}, \texttt{"X"}, \texttt{"Rx"}) to flat tuples or generator functions, with a debug assertion verifying unitarity $\|G^\dagger G - I\| < 10^{-12}$ on registration[cite: 1].

    \item A circuit model comprising \texttt{GateColumn}, \texttt{Circuit}, \texttt{Flatten}, and \texttt{Ascii}. \texttt{GateColumn} is a passive value class representing a time slice across qubit rows, with intra-column disjointness assertions preventing two operations from sharing target or control qubits. \texttt{Circuit} is an immutable list of columns with \texttt{withColumnAt}, \texttt{withWires}, and \texttt{numQubits}. \texttt{Flatten} implements a dual-layer pipeline: Layer 1 recursively inlines subcircuits into flat instruction objects carrying \texttt{virtualColumn} metadata; Layer 2 groups instructions sharing a virtual column into contiguous execution batches. \texttt{Ascii} ports Quirk's diagram parser to generate test circuits from text diagrams[cite: 1].

    \item A canonical JSON \texttt{Serializer} (\texttt{toJson}/\texttt{fromJson}) adhering to the v1 Canonical Schema: explicit \texttt{"format": "simqc-circuit"}, \texttt{"version": 1}, \texttt{"bitOrder": "little-endian"}, with control lines as explicit objects \texttt{controls: [{wire: 0, on: true}]}[cite: 1].

    \item A differential test harness (\texttt{kernel.test.js}, \texttt{Circuit.test.js}, \texttt{GateColumn.test.js}, \texttt{Serializer.test.js}) comparing the pairwise kernel against Sprint 1's $O(4^n)$ Kronecker reference over pseudo-random circuits for $n \le 6$, asserting norm preservation $\lVert\psi\rVert^2 \in 1.0 \pm 10^{-9}$ after each gate[cite: 1].
\end{enumerate}

The repository structure was finalized under \texttt{src/core}, \texttt{src/circuit}, \texttt{src/gates}, and \texttt{src/cli}, with \texttt{CONVENTIONS.md}, \texttt{FOLDER\_STRUCTURE.md}, and \texttt{SOURCE\_OF\_TRUTH.md} frozen as the source of truth[cite: 1].

\section{Learning: Matrix-Free Pairwise Simulation vs. Kronecker Products}
The textbook Kronecker product approach lifts a single-qubit gate $G$ on target qubit $t$ to the full $2^n \times 2^n$ operator:
\begin{equation}
    U = I \otimes \cdots \otimes G \otimes \cdots \otimes I,
\end{equation}
then multiplies $U$ against the state vector at $\mathcal{O}(4^n)$ cost[cite: 1]. At 16 qubits, a single lifted operator matrix requires $\sim$34~GB of RAM, and the matrix is overwhelmingly sparse---for $(I \otimes I \otimes H)$ on 3 qubits, 48 of the 64 entries are zero, meaning 75\% of the arithmetic operates on zero amplitudes[cite: 1].

The matrix-free pairwise method instead updates only the $2^{n-1}$ amplitude pairs whose indices differ in bit $t$ ($j = i \oplus 2^t$), each with the same $2\times2$ matrix $G$, yielding $\mathcal{O}(2^n)$ runtime per gate[cite: 1]. Our worked 3-qubit example confirmed that the pairwise method and the naive Kronecker product are numerically identical, and that CNOT reduces to pure in-place amplitude swapping (a permutation) requiring no floating-point multiplication[cite: 1].

\textbf{Entanglement forbids lazy Kronecker products.} We considered keeping each qubit as an independent $2\times1$ vector and deferring the tensor product to the end[cite: 1]. This holds only while the state is separable: the 2-qubit separability condition $c_{00}c_{11} - c_{01}c_{10} = 0$ holds after single-qubit gates but is violated the moment an entangling gate (CNOT, CZ) acts, since CNOT does not factorize as $A \otimes B$[cite: 1]. Simulators must therefore maintain the full $2^n$-dimensional vector from the start[cite: 1].

\textbf{Object allocation does not belong in the hot path.} Although Sprint 1's \texttt{Complex} class is complete and correct, representing the state vector as an array of \texttt{Complex} objects was rejected[cite: 1]. In V8, each object overhead costs roughly 40--56~bytes versus 16 bytes for two 64-bit floats in a \texttt{Float64Array} ($\approx 3\times$ memory reduction)[cite: 1]. Complex multiply-add with objects allocates temporary heap instances; at $2^{20}$ amplitudes per gate this triggers garbage-collector pauses that freeze the UI during drag-and-drop[cite: 1]. The kernel therefore operates exclusively on primitive doubles with zero heap allocation[cite: 1].

\section{Architecture Decisions}
\begin{itemize}
    \item \textbf{Layout:} little-endian indexing ($j = i \oplus (1 \ll q)$); two flat \texttt{Float64Array} buffers behind the \texttt{StateVector} facade[cite: 1].
    \item \textbf{Kernel arithmetic:} raw scalar doubles only; no \texttt{Complex} instances in gate-application loops[cite: 1].
    \item \textbf{Controlled gates:} pair update applied only when \texttt{(i \& controlMask) === controlMask}; non-matching amplitudes left untouched bit-for-bit[cite: 1].
    \item \textbf{Normalization:} norm verified after every gate; renormalize only if drift $> 10^{-9}$, since silent corrections mask arithmetic bugs[cite: 1].
    \item \textbf{Column semantics:} gates within one circuit column act on disjoint qubits; overlap prevented at construction[cite: 1].
    \item \textbf{Serialization:} explicit qubit indices; little-endian declared in schema metadata[cite: 1].
\end{itemize}

\section{Differential Validation}
The harness compares the $O(2^n)$ pairwise kernel against Sprint 1's $O(4^n)$ Kronecker reference over pseudo-random circuits for $n \le 6$, requiring agreement within $\epsilon \le 10^{-12}$ and norm $\lVert\psi\rVert^2 \in 1.0 \pm 10^{-9}$ after each gate[cite: 1]. Bidirectional CNOT tests ($0\to1$ swaps $\{1,3\},\{5,7\}$; $1\to0$ swaps $\{2,3\},\{6,7\}$) guard against endianness slips[cite: 1]. The Bell-state exit criterion ($H$ then CNOT$(0\to1)$ yields $(0.7071, 0, 0, 0.7071)$) is verified end-to-end through the headless CLI: \texttt{cat test/fixtures/bell\_state.json | node src/cli.js}[cite: 1]. Agreement between the naive Kronecker path and the pairwise kernel is maintained within tolerance across all tested depths and gate mixtures[cite: 1].

\section{Next Steps}
\begin{enumerate}
    \item Complete the headless CLI runner (\texttt{src/cli.js}) consuming canonical circuit JSON from \texttt{stdin} and outputting state-vector amplitudes or probabilities to \texttt{stdout}[cite: 1].
    \item Extend the differential harness to cover parameterized gates ($R_x, R_y, R_z$) and multi-qubit controlled operations[cite: 1].
    \item Begin Phase 1 of the AI module: local inference via Ollama/llama.cpp under the frozen \texttt{AI\_INTERFACE.md} contract, with \texttt{qwen2.5:3b} as the default development model[cite: 1].
    \item Explore multi-threading (Worker threads) for the kernel's block iteration to parallelize independent amplitude pairs[cite: 1].
\end{enumerate}

\textbf{Exit criterion:} differential validation passes for $n \le 6$ at $\epsilon \le 10^{-12}$, and the CLI produces Bell-state probabilities $[0.5, 0.0, 0.0, 0.5]$ from \texttt{bell\_state.json}[cite: 1].

\section{Risks and Mitigations}
\begin{enumerate}
    \item \textbf{Memory ceiling:} $\sim$256~MB at 24 qubits; mitigated by enforcing a hard simulation limit at 24 qubits with UI threshold warnings at 20 qubits[cite: 1].
    \item \textbf{GC pressure:} mitigated by the zero-allocation kernel design detailed in Section~II[cite: 1].
    \item \textbf{Bit-ordering inconsistencies:} mitigated by the frozen conventions specification and automated bidirectional differential tests against the reference implementation[cite: 1].
    \item \textbf{Numerical drift:} accumulated IEEE-754 rounding over deep circuits; mitigated by per-gate norm verification and selective renormalization at $10^{-9}$ drift threshold[cite: 1].
\end{enumerate}

\begin{thebibliography}{00}

\bibitem{quirk}
C.~Gidney, ``Quirk: A drag-and-drop quantum circuit simulator that runs in your browser,'' GitHub repository, Apache-2.0 license. [Online]. Available: \url{https://github.com/strilanc/quirk}

\bibitem{qiskit}
A.~Javadi-Abhari \emph{et~al.}, ``Qiskit: An open-source framework for quantum computing,'' IBM, 2017. [Online]. Available: \url{https://qiskit.org}

\bibitem{deraedt}
H.~De~Raedt and K.~Michielsen, ``Computational methods for simulating quantum computers,'' in \emph{Quantum Computing and Quantum Bits in Mesoscopic Systems}. Hoboken, NJ, USA: Wiley, 2004.

\bibitem{qhipster}
M.~Smelyanskiy, N.~P.~D.~Sawaya, and A.~Aspuru-Guzik, ``qHiPSTER: The Quantum High Performance Software Toolkit,'' 2016, \emph{arXiv:1601.07195}.

\bibitem{openqasm}
A.~W.~Cross \emph{et~al.}, ``OpenQASM 3 specification,'' 2021. [Online]. Available: \url{https://openqasm.com}

\end{thebibliography}

\end{document}
